\begin{exo}[Calcul d'espérance et de variance 1]
    \label{probas1}
    Soit $(X_p)_{p\in\N^*}$ une suite de variables aléatoires i.i.d suivant la loi uniforme sur $\inl1n$.
    On pose 
    \[
        N_n=\Card\lbrace X_1,\dots,X_n\rbrace  
    \]
    Calculer $\mathbb E(N_n)$ et $\mathbb V(N_n)$ et en donner des équivalents lorsque $n$ tend vers $+\infty$.
\end{exo}

\begin{solution}
    On va noter, pour $k\in\inl1n$, $S_k=\sum\mathbf1_{X_i=k}$, et alors
    \[
        N_n=\sum_{k=1}^n\mathbf1_{S_k\geq 1}    
    \]
    
    D'où, par linéarité de l'espérance 
    \[
        \mathbb E(N_n)=\sum_{k=1}^n\mathbb E(\mathbf1_{S_k\geq 1})=\sum_{k=1}^n\mathbb P(S_k\geq 1)
    \]

    Or
    \[
        \mathbb P(S_k\geq 1)=\mathbb P\left(\bigcup_{i=1}^n(X_i=k)\right)=1-\prod_{i=1}^n\mathbb P(X_i\neq k)    
    \]

    D'où $\mathbb P(S_k)=1-\left(\frac{n-1}n\right)^n$ pour tout $k$, soit $\mathbb E(N_n)=n\left[1-\left(\frac{n-1}n\right)^n\right]$.

    On exploite aussi la linéarité de l'espérance pour calculer la variance. 
   
    On a 
    \[
        \mathbb V(N_n)=\mathbb E(N_n^2)-\mathbb E(N_n)^2    
    \]

    Mais 
    \[
        N_n^2=\sum_{k=1}^n\mathbf1_{S_k\geq 1}^2+2\sum_{k<l}\mathbf1_{S_k\geq 1}\mathbf1_{S_l\geq 1}
    \]

    Donc $N_n^2=N_n+2\sum_{k<l}\mathbf1_{S_k\geq 1}\mathbf1_{S_l\geq 1}$ et 
    \[
        \mathbb E(N_n^2)=\mathbb E(N_n)+2\sum_{k<l}\mathbb E(\mathbf1_{S_k\geq 1}\mathbf1_{S_l\geq 1})
    \]
    Puis $\mathbb E(\mathbf1_{S_k\geq 1}\mathbf1_{S_l\geq 1})=\mathbb P(S_k \geq 1 \cap S_l \geq 1)$. 

    En passant par l'événement contraire, on obtient :
    \begin{align*}
        \mathbb P(S_k \geq 1 \cap S_l \geq 1) &= 1 - \mathbb P(S_k = 0 \cup S_l = 0)\\
        &= 1 - \mathbb P(S_k = 0) - \mathbb P(S_l = 0) + \mathbb P(S_k = 0 \cap S_l = 0)\\
        &= 1 - 2\left(\frac{n-1}{n}\right)^n + \left(\frac{n-2}{n}\right)^n
    \end{align*}

    Finalement, comme il y a $\frac{n(n-1)}{2}$ paires distinctes $(k,l)$ :
    \[
        \mathbb E(N_n^2) = \mathbb E(N_n) + n(n-1)\left[1 - 2\left(\frac{n-1}{n}\right)^n + \left(\frac{n-2}{n}\right)^n\right]
    \]

    Soit $\mathbb V(N_n) = \mathbb E(N_n^2) - \mathbb E(N_n)^2$.
    Donnons à présent des équivalents à ces quantités. 
    On a facilement $\mathbb E(N_n) \sim n(1-e^{-1})$. Pour la variance, un développement asymptotique plus fin (ordre $1/n$) des termes $\left(1 - \frac{x}{n}\right)^n$ montre que les termes en $n^2$ s'annulent exactement, laissant une variance linéaire en $n$ :
    \[
        \mathbb V(N_n) \sim n(e^{-1} - 2e^{-2})
    \]
\end{solution}

\begin{exo}[Calcul d'espérance et de variance 2]
    \label{probas2}
    On note $N_n$ le nombre de points fixes d'une permutation aléatoire suivant la loi uniforme sur $\mathfrak S_n$.
    Calculer l'espérance et la variance de $N_n$.
\end{exo}

\begin{solution}
    Comme dans l'exercice précédent, on va exploiter la linéarité de l'espérance. 
    On prend $S$ une permutation aléatoire et on écrit 
    \[
        N_n=\sum_{k=1}^n\mathbf1_{S(k)=k}
    \]
    Ainsi, par linéarité de l'espérance
    \[
        \mathbb E(N_n)=\sum_{k=1}^n \mathbb P(S(k)=k)
    \]
    Maintenant, à $k$ fixé, que vaut $\mathbb P(S(k)=k)$ ?
    Comme $S$ suit une loi uniforme, on cherche alors le nombre de permutations dans $\mathfrak S_n$ fixant $k$. 
    On considère l'application $\theta i\in\inl1n\mapsto \mathbf1_{i<k}i+\mathbf1_{i > k}(i-1)$, et on pose 
    $\Gamma : \sigma \in \mathfrak S_{n-1}\mapsto [l\in\inl1n\mapsto \mathbf1_{l\neq k}\sigma\circ\theta(l)+\mathbf1_{l=k}k]\in\mathfrak S_n$.
    On vérifie facilement que cette application réalise une bijection entre l'ensemble des permutations de $\mathfrak S_n$ fixant et $k$ et les permutations de $\mathfrak S_{n-1}$
    (pour $n\geq 2$) ; on en déduit $\mathbb P(S(k)=k)=1/n$, puis $\mathbb E(N_n)=1$.

    Pour la variance, on a 
    \[
        \mathbb V(N_n)=\mathbb E(N_n^2)-\mathbb E(N_n)^2=\mathbb E(N_n^2)-1
    \]
    Mais 
    \[
        N_n^2=\sum_{k=1}^n\mathbf1_{S(k)=k}+2\sum_{k<l}\mathbf1_{S(k)=k}\mathbf1_{S(l)=l}
    \]
    Donc $\mathbb E(N_n^2)=\mathbb E(N_n)+2\sum_{k<l}\mathbb E(\mathbf1_{S(k)=k}\mathbf1_{S(l)=l})$.
    Or pour $k<l$, $\mathbb E(\mathbf1_{S(k)=k}\mathbf1_{S(l)=l})=\mathbb P([S(k)=k]\cap[S(l)=l])$.
    Comme plus haut, on cherche cette fois-ci à compter les permutations qui fixent deux points ; on laisse 
    au lecteur les sois de vérifier qu'elles sont au nombre de $(n-2)!$, et alors 
    \[
        \mathbb E(\mathbf1_{S(k)=k}\mathbf1_{S(l)=l})=\frac1{n(n-1)}
    \]
    Soit $\mathbb E(N_n^2)=\mathbb E(N_n)+2\frac{n(n-1)}2\times\frac1{n(n-1)}=1+1=2$.
    Finalement, $\mathbb V(N_n)=\mathbb E(N_n^2)-1=2-1=1$.

\end{solution}